\documentclass[11pt]{article}
\usepackage[T1]{fontenc}
\usepackage[utf8]{inputenc}
\usepackage{lmodern}
\usepackage{microtype}
\usepackage[margin=1.05in]{geometry}
\usepackage{amsmath,amssymb,mathtools,amsthm}
% Canonical mathematical notation for active FabiusFunction LaTeX documents.
%
% This file is intentionally a plain TeX input rather than a LaTeX package.
% Every standalone document loads it by an explicit path relative to that
% document's directory; included fragments inherit the definitions from their
% root document.  Keep this file independent of document layout and metadata.
%
% Contract:
%   * one command has one mathematical meaning throughout the active corpus;
%   * names encode semantic roles rather than a document-local choice of letter;
%   * visually similar but differently normalized objects have different print;
%   * every command below is defined normatively in the notation catalogue.
%
% The active corpus excludes docs/archive/.  Historical sources may retain
% their original notation only inside an explicitly labelled quotation.

\makeatletter
\@ifundefined{mathbb}{%
  \PackageError{fabius-notation}{The amssymb package must be loaded first}%
    {Load amsmath and amssymb before inputting fabius-notation.tex.}%
}{}
\@ifundefined{operatorname}{%
  \PackageError{fabius-notation}{The amsmath package must be loaded first}%
    {Load amsmath before inputting fabius-notation.tex.}%
}{}
\makeatother

% Version stamp used by the catalogue and by static audits.
\newcommand{\FabiusNotationVersion}{2026-08-31}

% Number systems.  NaturalNumbers is explicitly the nonnegative convention.
\newcommand{\NaturalNumbers}{\mathbb{N}_{0}}
\newcommand{\PositiveIntegers}{\mathbb{N}_{>0}}
\newcommand{\IntegerNumbers}{\mathbb{Z}}
\newcommand{\RationalNumbers}{\mathbb{Q}}
\newcommand{\RealNumbers}{\mathbb{R}}
\newcommand{\ComplexNumbers}{\mathbb{C}}
\newcommand{\ExtendedRealNumbers}{\overline{\mathbb{R}}}

% The three principal Fabius--Rvachev functions.  Subscripts are deliberate:
% a bare F and a calligraphic F have several unrelated uses in the corpus.
\newcommand{\FabiusBounded}{F_{\mathrm{bd}}}
\newcommand{\FabiusGlobal}{F_{\mathrm{gl}}}
\newcommand{\FabiusQuantile}{Q_{F}}
\newcommand{\FabiusClampedQuantile}{Q_{F,\mathrm{cl}}}
\newcommand{\RvachevUp}{\operatorname{up}}

% Constants, elementary operators, and delimiters.
\newcommand{\EulerE}{\mathrm{e}}
\newcommand{\ImaginaryUnit}{\mathrm{i}}
\newcommand{\Differential}{\mathop{}\!\mathrm{d}}
\newcommand{\AbsoluteValue}[1]{\left\lvert #1\right\rvert}
\newcommand{\Norm}[1]{\left\lVert #1\right\rVert}
\newcommand{\Floor}[1]{\left\lfloor #1\right\rfloor}
\newcommand{\Ceiling}[1]{\left\lceil #1\right\rceil}
\newcommand{\FractionalPart}[1]{\left\{#1\right\}}
\newcommand{\PositivePart}[1]{\left(#1\right)_{+}}
\newcommand{\IndicatorOf}[1]{\mathbf{1}_{#1}}
\newcommand{\SetOf}[1]{\left\{#1\right\}}
\newcommand{\InnerProduct}[1]{\left\langle #1\right\rangle}
\newcommand{\CardinalityOf}[1]{\left\lvert #1\right\rvert}
\newcommand{\DefinitionEquals}{\mathrel{\coloneqq}}
\newcommand{\Sign}{\operatorname{sgn}}
\newcommand{\IdentityOperator}{\operatorname{id}}
\newcommand{\SpanOperator}{\operatorname{span}}
\newcommand{\TraceOperator}{\operatorname{tr}}
\newcommand{\RankOperator}{\operatorname{rank}}
\newcommand{\DiagonalOperator}{\operatorname{diag}}
\newcommand{\ResidueOperator}{\operatorname*{Res}}
\newcommand{\LeastCommonMultiple}{\operatorname{lcm}}
\newcommand{\OrderOperator}{\operatorname{ord}}
\newcommand{\PrincipalLogarithm}{\operatorname{Log}}
\newcommand{\FallingFactorial}[2]{\left(#1\right)^{\underline{#2}}}
\newcommand{\RisingFactorial}[2]{\left(#1\right)^{\overline{#2}}}

% Support, topology, convolution, and metric notation.
\newcommand{\SupportOperator}{\operatorname{supp}}
\newcommand{\SupportOf}[1]{\SupportOperator\!\left(#1\right)}
\newcommand{\TopologicalSupportOf}[1]{\operatorname{tsupp}\!\left(#1\right)}
\newcommand{\Convolution}{\mathbin{*}}
\newcommand{\MetricDistance}{\operatorname{dist}}
\newcommand{\TotalVariationDistance}{d_{\mathrm{TV}}}

% Probability.  Equality and convergence in law are not metric distance.
\newcommand{\Probability}{\mathbb{P}}
\newcommand{\Expectation}{\mathbb{E}}
\newcommand{\Variance}{\operatorname{Var}}
\newcommand{\Covariance}{\operatorname{Cov}}
\newcommand{\LawOperator}{\operatorname{Law}}
\newcommand{\LawOf}[1]{\LawOperator\!\left(#1\right)}
\newcommand{\EqualInLaw}{\mathrel{\overset{\mathrm{law}}{=}}}
\newcommand{\ConvergesInLaw}{\mathrel{\xrightarrow{\mathrm{law}}}}

% Fourier and integral transforms.  The printed labels expose normalization.
% FourierTwoPi uses exp(-2*pi*i*x*xi); FourierAngular uses exp(-i*x*omega).
\newcommand{\FourierTwoPi}[1]{\widehat{#1}_{\,2\pi}}
\newcommand{\FourierAngular}[1]{\widehat{#1}_{\,\mathrm{ang}}}
\newcommand{\LaplaceTransformSymbol}{\mathcal{L}}
\newcommand{\MellinTransformSymbol}{\mathcal{M}}
\newcommand{\LaplaceTransformOf}[1]{\LaplaceTransformSymbol\!\left[#1\right]}
\newcommand{\MellinTransformOf}[1]{\MellinTransformSymbol\!\left[#1\right]}
\newcommand{\MomentGeneratingFunction}[1]{M_{#1}}
\newcommand{\CharacteristicFunction}[1]{\varphi_{#1}}

% The two standard sinc functions must never share one symbol.
\newcommand{\SincRad}{\operatorname{sinc}_{\mathrm{rad}}}
\newcommand{\SincPi}{\operatorname{sinc}_{\pi}}

% Binary arithmetic and Thue--Morse notation.
\newcommand{\BinaryDigitSum}{\operatorname{wt}_{2}}
\newcommand{\ThueMorseSignSymbol}{\varepsilon}
\newcommand{\ThueMorseSign}[1]{\varepsilon_{#1}}
\newcommand{\PadicValuation}[1]{\nu_{#1}}
\newcommand{\TwoAdicValuation}{\nu_{2}}

% q-series notation.  Every base is an explicit argument.
\newcommand{\QPochhammer}[3]{\left(#1;#2\right)_{#3}}
\newcommand{\GaussianBinomial}[3]{%
  \genfrac{[}{]}{0pt}{}{#1}{#2}_{#3}%
}
\newcommand{\QInteger}[2]{\left[#1\right]_{#2}}

% Named special functions.
\newcommand{\Polylogarithm}{\operatorname{Li}}
\newcommand{\LambertW}{W}
\newcommand{\GammaFunction}{\Gamma}
\newcommand{\RiemannZeta}{\zeta}

% Asymptotics.  BigO and LittleO are operator tokens for existing formula
% grammar; prefer the At forms so that the limiting regime is printed.
\newcommand{\BigO}{\mathcal{O}}
\newcommand{\LittleO}{o}
\newcommand{\BigOAt}[2]{\mathcal{O}_{#2}\!\left(#1\right)}
\newcommand{\LittleOAt}[2]{o_{#2}\!\left(#1\right)}

\usepackage{booktabs,longtable,array,tabularx}
\usepackage{enumitem}
\usepackage{xcolor}
\usepackage{hyperref}
\usepackage[nameinlink,capitalise,noabbrev]{cleveref}
\usepackage{fancyhdr}
\usepackage{bookmark}
\usepackage{mathrsfs}
\usepackage{bm}
\usepackage{verbatim}

\hypersetup{
  colorlinks=true,
  linkcolor=blue!45!black,
  citecolor=green!35!black,
  urlcolor=blue!55!black,
  pdftitle={A Proof-Driven Guide to q-Series and Rogers--Ramanujan Theory},
  pdfauthor={OpenAI, prepared from the supplied q-series archive}
}
\setlength{\parindent}{1.25em}
\setlength{\parskip}{0.25em}
\allowdisplaybreaks
\numberwithin{equation}{section}
\pagestyle{fancy}
\fancyhf{}
\fancyhead[L]{\small A Proof-Driven Guide to $q$-Series}
\fancyhead[R]{\small\thepage}
\renewcommand{\headrulewidth}{0.3pt}

\newtheorem{theorem}{Theorem}[section]
\newtheorem{lemma}[theorem]{Lemma}
\newtheorem{proposition}[theorem]{Proposition}
\newtheorem{corollary}[theorem]{Corollary}
\newtheorem{identity}[theorem]{Identity}
\newtheorem{conjecture}[theorem]{Conjecture}
\theoremstyle{definition}
\newtheorem{definition}[theorem]{Definition}
\newtheorem{example}[theorem]{Example}
\newtheorem{problem}[theorem]{Problem}
\theoremstyle{remark}
\newtheorem{remark}[theorem]{Remark}

\newcommand{\CT}{\operatorname{CT}}

\title{\textbf{A Proof-Driven Guide to $q$-Series,\\Basic Hypergeometric Identities, Bailey Chains,\\and Rogers--Ramanujan Theory}\\[0.7em]
\large A unified article based on the supplied $q$-series archive and classical primary literature}
\author{Prepared by OpenAI}
\date{August 2026}

\begin{document}
\maketitle

\begin{abstract}
The supplied archive contains a broad collection of notes and formula sheets on
$q$-Pochhammer symbols, Gaussian coefficients, basic hypergeometric series,
Rogers--Ramanujan type identities, Bailey pairs, theta and eta products,
partition generating functions, modulus-specific identities, and
Borwein-type sign phenomena.  This article reorganizes that material into one
proof-driven development.  Every result labelled as a theorem, proposition,
lemma, corollary, or identity is proved here; statements not proved are
explicitly labelled as conjectures or research problems.

The main line begins with the finite and infinite $q$-binomial theorems and
continues through Heine's transformation, the $q$-Gauss and terminating
$q$-Vandermonde sums, a finite Rogers product lemma, and the
$q$-Pfaff--Saalsch\"utz summation.  These results are then used to prove
Bailey inversion, Bailey's lemma, its limiting forms, and a parameter-lowering
step.  From this machinery we obtain complete analytic proofs of both
Rogers--Ramanujan identities and the entire Andrews--Gordon family.  A separate
argument derives the Rogers--Ramanujan continued fraction from finite
continuants.  Jacobi's triple product, the quintuple product, Poisson summation,
theta transformations, the Dedekind eta transformation, and sharp
$q\to1^-$ asymptotics are proved as well.  The final sections explain how the
moduli $7,9,14,$ and $27$ material in the archive fits into the general theory,
and formulate carefully separated open directions involving finite analogues,
Bailey lattices, parity restrictions, and coefficient-sign problems.
\end{abstract}

\tableofcontents
\newpage

\section{How the source material is organized here}

The archive mixes several conventions: some files use $q$ for the nome
$\EulerE^{\pi i\tau}$, others use $q=\EulerE^{2\pi i\tau}$; some suppress the base in
$(a;q)_n$; and the normalization of ${}_r\phi_s$ is not uniform.  We use one
convention throughout and translate all identities into it before proving
them.  The resulting dependency graph is
\[
\begin{gathered}
\text{finite $q$-binomial theorem}
 \Longrightarrow \text{infinite $q$-binomial theorem}
 \Longrightarrow \text{Heine and $q$-Gauss},\\
\text{finite Cauchy/Rogers identities}
 \Longrightarrow q\text{-Pfaff--Saalsch\"utz}
 \Longrightarrow \text{Bailey lemma},\\
\text{Bailey chains and a parameter shift}
 \Longrightarrow \text{Rogers--Ramanujan and Andrews--Gordon},\\
\text{Euler expansions and Laurent functional equations}
 \Longrightarrow \text{triple/quintuple products},\\
\text{Poisson summation plus products}
 \Longrightarrow \text{theta/eta transformations and asymptotics}.
\end{gathered}
\]
The modulus-specific files are treated as specializations or extensions of
this graph instead of as disconnected lists of formulas.

\subsection*{Logical and analytic convention}
Most finite identities are polynomial or rational identities and therefore
hold formally whenever no displayed denominator vanishes.  Infinite products
and series are first proved for $|q|<1$ in their natural domains of absolute
convergence.  Formal-power-series interpretations are stated separately.
When a parameter value is singular in an intermediate formula but the final
expression has a removable singularity, the value is understood by a limit.

\section{$q$-shifted factorials and Gaussian coefficients}

\begin{definition}[$q$-shifted factorial]
For $n\in\NaturalNumbers$ and a base $q$, put
\[
(a;q)_0=1,\qquad (a;q)_n=\prod_{j=0}^{n-1}(1-aq^j),
\]
and, when $|q|<1$,
\[
(a;q)_\infty=\prod_{j=0}^{\infty}(1-aq^j).
\]
For a list of parameters,
\[
(a_1,\ldots,a_r;q)_n=\prod_{j=1}^r(a_j;q)_n.
\]
\end{definition}

The following elementary rules will be used without repeated comment:
\begin{align}
(a;q)_{m+n}&=(a;q)_m(aq^m;q)_n,\label{eq:split}\
(aq^m;q)_n&=\frac{(a;q)_{m+n}}{(a;q)_m},\label{eq:shift}\
(a;q)_n&=(-a)^nq^{\binom n2}(q^{1-n}/a;q)_n.\label{eq:reverse}
\end{align}
All three follow directly by rearranging their finite products.

\begin{lemma}[Convergence of the infinite product]\label{lem:prodconv}
Let $|q|<1$.  The product $(a;q)_\infty$ converges locally uniformly as a
function of $a$, hence defines an entire function of $a$.  It vanishes exactly
at $a=q^{-m}$, $m\in\NaturalNumbers$, with simple zeros when $q\ne0$.
\end{lemma}
\begin{proof}
On a compact set $|a|\le M$,
\(\sum_{j\ge0}|aq^j|\le M/(1-|q|)\).  The standard convergence criterion for
products $\prod(1+u_j)$, applied after discarding finitely many factors if
necessary, gives uniform convergence.  Uniform limits of polynomials are
holomorphic.  A factor vanishes precisely when $a=q^{-m}$; all remaining
factors are nonzero, and the vanishing factor is linear in $a$.
\end{proof}

\begin{definition}[Gaussian coefficient]
For $0\le k\le n$ define
\[
\GaussianBinomial{n}{k}{q}=\frac{(q;q)_n}{(q;q)_k(q;q)_{n-k}},
\]
and set it equal to $0$ outside this range.
\end{definition}

\begin{proposition}[Pascal recurrences]\label{prop:qpascal}
For $n\ge0$ and all integers $k$,
\begin{align}
\GaussianBinomial{n+1}{k}{q}&=\GaussianBinomial{n}{k}{q}+q^{n+1-k}\GaussianBinomial{n}{k-1}{q},\label{eq:pascal1}\\
\GaussianBinomial{n+1}{k}{q}&=q^k\GaussianBinomial{n}{k}{q}+\GaussianBinomial{n}{k-1}{q}.\label{eq:pascal2}
\end{align}
\end{proposition}
\begin{proof}
After putting the two terms over the common denominator
$(q;q)_k(q;q)_{n+1-k}$, the numerator in \eqref{eq:pascal1} is
\[
(q;q)_n\{1-q^{n+1-k}+q^{n+1-k}(1-q^k)\}
=(q;q)_n(1-q^{n+1}).
\]
This is $(q;q)_{n+1}$.  The second recurrence follows in the same way, or from
the first after the symmetry $\GaussianBinomial{n}{k}{q}=\GaussianBinomial{n}{n-k}{q}$.
\end{proof}

\begin{remark}[Combinatorial meaning]
When $q$ is a prime power, $\GaussianBinomial{n}{k}{q}$ counts the $k$-dimensional subspaces of
$\mathbb F_q^n$.  As a polynomial in an indeterminate $q$, it is also the
rank-generating function of partitions whose Ferrers diagrams fit inside a
$k\times(n-k)$ rectangle.  Neither interpretation is needed for the analytic
proofs below, but both explain why positivity and convolution identities are
natural.
\end{remark}

\section{The finite and infinite $q$-binomial theorems}

\begin{theorem}[Finite $q$-binomial theorem]\label{thm:finiteqbt}
For every $n\in\NaturalNumbers$,
\begin{equation}\label{eq:finiteqbt}
(z;q)_n=\sum_{k=0}^n(-1)^kq^{\binom k2}\GaussianBinomial{n}{k}{q} z^k.
\end{equation}
\end{theorem}
\begin{proof}
The assertion is clear for $n=0$.  If it holds for $n$, multiply by
$1-zq^n$.  The coefficient of
$(-1)^kq^{\binom k2}z^k$ in the resulting expression is
\[
\GaussianBinomial{n}{k}{q}+q^{n+1-k}\GaussianBinomial{n}{k-1}{q}=\GaussianBinomial{n+1}{k}{q}
\]
by \eqref{eq:pascal1}.  This proves the induction step.
\end{proof}

\begin{corollary}[Finite subset sums]\label{cor:subsets}
For $0\le k\le n$,
\[
\sum_{0\le j_1<\cdots<j_k<n}q^{j_1+\cdots+j_k}
=q^{\binom k2}\GaussianBinomial{n}{k}{q}.
\]
\end{corollary}
\begin{proof}
Compare the coefficient of $(-z)^k$ in the product
$\prod_{j=0}^{n-1}(1-zq^j)$ and in \eqref{eq:finiteqbt}.
\end{proof}

\begin{theorem}[Infinite $q$-binomial theorem]\label{thm:qbt}
If $|q|<1$ and $|z|<1$, then
\begin{equation}\label{eq:qbt}
\sum_{n=0}^{\infty}\frac{(a;q)_n}{(q;q)_n}z^n
=\frac{(az;q)_\infty}{(z;q)_\infty}.
\end{equation}
The identity also holds coefficientwise in $\ComplexNumbers[a,q][[z]]$.
\end{theorem}
\begin{proof}
Let
\(F(z)=\sum_{n\ge0}c_nz^n\), where
$c_n=(a;q)_n/(q;q)_n$.  The ratio
\[
\frac{c_n}{c_{n-1}}=\frac{1-aq^{n-1}}{1-q^n}
\]
is equivalent, coefficient by coefficient, to
\begin{equation}\label{eq:qbtfe}
(1-z)F(z)=(1-az)F(qz),\qquad F(0)=1.
\end{equation}
The product
$P(z)=(az;q)_\infty/(z;q)_\infty$ satisfies the same equation, because
\[
\frac{P(z)}{P(qz)}=\frac{1-az}{1-z},
\]
and $P(0)=1$.  Equation \eqref{eq:qbtfe} recursively determines every Taylor
coefficient, so $F=P$ formally.  For $|z|<1$, the series converges absolutely
by the ratio test and the products converge locally uniformly by
\cref{lem:prodconv}; hence the formal identity is also analytic.
\end{proof}

\begin{corollary}[Euler's two expansions]\label{cor:eulerexp}
For $|q|<1$,
\begin{align}
\frac1{(z;q)_\infty}&=\sum_{n=0}^{\infty}\frac{z^n}{(q;q)_n},
&& |z|<1,\label{eq:euler1}\\
(z;q)_\infty&=\sum_{n=0}^{\infty}
\frac{(-1)^nq^{\binom n2}}{(q;q)_n}z^n.
\label{eq:euler2}
\end{align}
\end{corollary}
\begin{proof}
Set $a=0$ in \eqref{eq:qbt} to obtain \eqref{eq:euler1}.  In
\eqref{eq:qbt}, replace $z$ by $z/a$ and let $|a|\to\infty$.  For each fixed
$n$,
\[
(a;q)_n(z/a)^n\longrightarrow(-1)^nq^{\binom n2}z^n,
\]
while the product side tends to $(z;q)_\infty$.  On compact $z$-sets the
series are dominated by a convergent geometric-majorant refinement, so the
limit may be taken termwise.  The same argument is valid coefficientwise.
\end{proof}

\begin{corollary}[Reciprocal finite product]\label{cor:reciprocalfinite}
For an integer $m\ge0$ and $|z|<1$,
\[
\frac1{(z;q)_{m+1}}=
\sum_{n=0}^{\infty}\GaussianBinomial{m+n}{n}{q}z^n.
\]
\end{corollary}
\begin{proof}
Apply \eqref{eq:qbt} with $a=q^{m+1}$.  The product quotient is
$1/(z;q)_{m+1}$, and
\[
\frac{(q^{m+1};q)_n}{(q;q)_n}
=\frac{(q;q)_{m+n}}{(q;q)_m(q;q)_n}=\GaussianBinomial{m+n}{n}{q}.
\]
\end{proof}

\begin{lemma}[Finite Cauchy augmentation]\label{lem:cauchyaug}
For $n\ge0$,
\begin{equation}\label{eq:cauchyaug}
(xy;q)_n=\sum_{k=0}^n\GaussianBinomial{n}{k}{q}(x;q)_k y^k(y;q)_{n-k}.
\end{equation}
\end{lemma}
\begin{proof}
For $n=0$ this is immediate.  Denote the right-hand side by $C_n(x,y)$.
Use \eqref{eq:pascal1}, split the sum, shift $k\mapsto k+1$ in the second
part, and use
$(x;q)_{k+1}=(x;q)_k(1-xq^k)$ and
$(y;q)_{n+1-k}=(y;q)_{n-k}(1-yq^{n-k})$.
The two parts combine to
\[
C_{n+1}(x,y)=(1-xyq^n)C_n(x,y).
\]
Since $C_0=1$, induction gives $C_n=(xy;q)_n$.
\end{proof}

\section{Basic hypergeometric series and Heine's transformation}

\begin{definition}[Basic hypergeometric series]
For parameters away from poles, define
\[
{}_r\phi_s\!\left(\begin{matrix}a_1,\dots,a_r\\b_1,\dots,b_s\end{matrix};q,z\right)
=\sum_{n=0}^{\infty}
\frac{(a_1,\dots,a_r;q)_n}{(q,b_1,\dots,b_s;q)_n}
\bigl((-1)^nq^{\binom n2}\bigr)^{1+s-r}z^n.
\]
When one upper parameter is $q^{-N}$, the series terminates at $n=N$.
\end{definition}

\begin{theorem}[Heine's first transformation]\label{thm:heine}
Assume $|q|<1$, $|z|<1$, and initially $|b|<1$.  Then
\begin{equation}\label{eq:heine}
{}_2\phi_1\!\left(\begin{matrix}a,b\\c\end{matrix};q,z\right)
=\frac{(b,az;q)_\infty}{(c,z;q)_\infty}
{}_2\phi_1\!\left(\begin{matrix}c/b,z\\az\end{matrix};q,b\right).
\end{equation}
Both sides extend meromorphically to all parameters for which the displayed
expressions are defined.
\end{theorem}
\begin{proof}
Use
\[
\frac{(b;q)_n}{(c;q)_n}
=\frac{(b;q)_\infty}{(c;q)_\infty}
 \frac{(cq^n;q)_\infty}{(bq^n;q)_\infty}
\]
and expand the last quotient by \cref{thm:qbt}:
\[
\frac{(cq^n;q)_\infty}{(bq^n;q)_\infty}
=\sum_{k=0}^{\infty}\frac{(c/b;q)_k}{(q;q)_k}(bq^n)^k.
\]
Absolute convergence permits interchange of the $n$- and $k$-sums.  The
inner $n$-sum is another $q$-binomial theorem:
\[
\sum_{n\ge0}\frac{(a;q)_n}{(q;q)_n}(zq^k)^n
=\frac{(azq^k;q)_\infty}{(zq^k;q)_\infty}
=\frac{(az;q)_\infty}{(z;q)_\infty}
  \frac{(z;q)_k}{(az;q)_k}.
\]
Substitution gives exactly the right-hand side of \eqref{eq:heine}.  The
identity theorem supplies meromorphic continuation.
\end{proof}

\begin{corollary}[$q$-Gauss summation]\label{cor:qgauss}
If $|c/(ab)|<1$, then
\begin{equation}\label{eq:qgauss}
{}_2\phi_1\!\left(\begin{matrix}a,b\\c\end{matrix};q,\frac{c}{ab}\right)
=\frac{(c/a,c/b;q)_\infty}{(c,c/(ab);q)_\infty}.
\end{equation}
\end{corollary}
\begin{proof}
Set $z=c/(ab)$ in \eqref{eq:heine}.  Since $az=c/b$, one upper and one lower
parameter in the transformed series cancel.  The remaining series is
\[
\sum_{k\ge0}\frac{(c/(ab);q)_k}{(q;q)_k}b^k
=\frac{(c/a;q)_\infty}{(b;q)_\infty}
\]
by \cref{thm:qbt}.  Multiplying by Heine's prefactor yields
\eqref{eq:qgauss}.
\end{proof}

\begin{corollary}[First terminating $q$-Vandermonde sum]\label{cor:qvand1}
For $N\in\NaturalNumbers$,
\begin{equation}\label{eq:qvand1}
{}_2\phi_1\!\left(\begin{matrix}q^{-N},b\\c\end{matrix};q,
\frac{cq^N}{b}\right)
=\frac{(c/b;q)_N}{(c;q)_N}.
\end{equation}
\end{corollary}
\begin{proof}
Set $a=q^{-N}$ in \eqref{eq:qgauss}.  The series terminates, so no
convergence restriction on its argument remains.  On the product side,
\[
\frac{(cq^N;q)_\infty}{(c;q)_\infty}=\frac1{(c;q)_N},\qquad
\frac{(c/b;q)_\infty}{(cq^N/b;q)_\infty}=(c/b;q)_N.
\]
\end{proof}

\begin{proposition}[Second terminating $q$-Vandermonde sum]\label{prop:qvand2}
For $N\in\NaturalNumbers$,
\begin{equation}\label{eq:qvand2}
{}_2\phi_1\!\left(\begin{matrix}q^{-N},b\\c\end{matrix};q,q\right)
=b^N\frac{(c/b;q)_N}{(c;q)_N}.
\end{equation}
\end{proposition}
\begin{proof}
Multiply both sides by $(c;q)_N$ and regard them as polynomials in $b$ of
degree at most $N$.  It is enough to verify the equality at the $N+1$ points
$b=cq^j$, $0\le j\le N$.  At such a point, split
$(cq^j;q)_k/(c;q)_k=(cq^k;q)_j/(c;q)_j$ using \eqref{eq:split}; then use the
finite $q$-binomial theorem on the resulting sum.  Explicitly,
\begin{align*}
(c;q)_N S_N(cq^j)
&=\sum_{k=0}^N
 \frac{(q^{-N};q)_k}{(q;q)_k}q^k
 (cq^k;q)_j(cq^k;q)_{N-k}\\
&=(c;q)_j\sum_{k=0}^N
 \frac{(q^{-N};q)_k}{(q;q)_k}q^k
 (cq^{j+k};q)_{N-j-k},
\end{align*}
where terms with a negative final length are interpreted as zero after the
finite cancellation that produced the line.  Reversing the finite product
and applying \eqref{eq:finiteqbt} gives
\[
(cq^j)^N(q^{-j};q)_N,
\]
which equals $(cq^j)^N(q^{-j};q)_N=(cq^j)^N(c/(cq^j);q)_N$ and therefore is
the claimed right-hand side at $b=cq^j$.  Polynomial interpolation completes
the proof.  (The displayed manipulation is also the usual polynomial proof
of the second $q$-Vandermonde formula.)
\end{proof}

\section{A finite Rogers product lemma and $q$-Pfaff--Saalsch\"utz}

\begin{lemma}[Rogers finite product lemma]\label{lem:rogersfinite}
For $n\in\NaturalNumbers$,
\begin{equation}\label{eq:rogersfinite}
(at;q)_n(bt;q)_n
=\sum_{k=0}^n\GaussianBinomial{n}{k}{q}(a;q)_k(b;q)_k t^k
(t;q)_{n-k}(abtq^k;q)_{n-k}.
\end{equation}
\end{lemma}
\begin{proof}
Put
\[
S_{n,k}=\GaussianBinomial{n}{k}{q}(a;q)_k(b;q)_k t^k
(t;q)_{n-k}(abtq^k;q)_{n-k},
\qquad R_n=\sum_{k=0}^n S_{n,k}.
\]
We prove by induction that
$R_n=(at;q)_n(bt;q)_n$.  The case $n=0$ is immediate.  Set
$Q=q^n$ and, in the $k$th summand, $u=q^k$.  Pascal's recurrence
\eqref{eq:pascal1}, followed only by the splitting rule
\eqref{eq:split}, gives
\begin{equation}\label{eq:Rogers-rec-expansion}
R_{n+1}=(1-abtQ)\sum_{k=0}^n S_{n,k}E(u),
\qquad
E(u)=1-\frac{tQ}{u}
+\frac{tQ}{u}\frac{(1-au)(1-bu)}{1-abtu}.
\end{equation}
The required multiplier is
$M=(1-atQ)(1-btQ)$.  The following rational certificate is obtained by
ordinary expansion and factorization:
\[
\mathcal C(u)=\frac{Q a b t \left(u - 1\right) \left(Q t - u\right)}{u}.
\]
It satisfies the exact rational identity
\begin{equation}\label{eq:Rogers-certificate}
\mathcal C(qu)\frac{S_{n,k+1}}{S_{n,k}}-\mathcal C(u)
=(1-abtQ)E(u)-M,
\end{equation}
where
\[
\frac{S_{n,k+1}}{S_{n,k}}
=\frac{t(1-Q/u)(1-au)(1-bu)}
{(1-qu)(1-tQ/(qu))(1-abtu)}.
\]
Equation \eqref{eq:Rogers-certificate} is verified by multiplying by the
three displayed denominators; the two sides reduce to the same polynomial in
$u$.  Multiplying it by $S_{n,k}$ and summing from $k=0$ to $n$ telescopes:
\[
\sum_{k=0}^n S_{n,k}\bigl((1-abtQ)E(q^k)-M\bigr)
=\sum_{k=0}^n
\bigl(\mathcal C(q^{k+1})S_{n,k+1}
      -\mathcal C(q^k)S_{n,k}\bigr)=0.
\]
The lower boundary vanishes because $\mathcal C(1)=0$; the upper one
vanishes because $S_{n,n+1}=0$, equivalently because the ratio contains
$1-Q/u$ at $u=Q$.  Thus \eqref{eq:Rogers-rec-expansion} gives
$R_{n+1}=MR_n$, and induction yields
$R_n=(at;q)_n(bt;q)_n$.
\end{proof}

\begin{theorem}[$q$-Pfaff--Saalsch\"utz summation]\label{thm:qps}
For $n\in\NaturalNumbers$,
\begin{equation}\label{eq:qps}
{}_3\phi_2\!\left(
\begin{matrix}a,b,q^{-n}\\c,abq^{1-n}/c\end{matrix};q,q\right)
=\frac{(c/a,c/b;q)_n}{(c,c/(ab);q)_n}.
\end{equation}
\end{theorem}
\begin{proof}
Multiply \eqref{eq:qps} by $(c;q)_n(c/(ab);q)_n$.  In its $k$th summand use
\eqref{eq:reverse} twice, once for $(q^{-n};q)_k$ and once for
$(abq^{1-n}/c;q)_k$.  All signs and powers of $q$ cancel, and the summand
becomes
\[
\GaussianBinomial{n}{k}{q}(a;q)_k(b;q)_k\left(\frac{c}{ab}\right)^k
(c/(ab);q)_{n-k}(cq^k;q)_{n-k}.
\]
Consequently the denominator-cleared identity is exactly
\cref{lem:rogersfinite} with $t=c/(ab)$.  Its left-hand side is
$(c/b;q)_n(c/a;q)_n$, which is the denominator-cleared right-hand side of
\eqref{eq:qps}.
\end{proof}

\begin{remark}
The preceding proof is useful structurally: Bailey's lemma rests on a wholly
finite polynomial identity, not on an appeal to analytic continuation or an
unproved summation formula.  This was one of the main proof-repair goals in
reorganizing the source archive.
\end{remark}

\section{Jacobi's triple product and the quintuple product}

\begin{theorem}[Jacobi triple product]\label{thm:jtp}
For $|q|<1$ and $z\ne0$,
\begin{equation}\label{eq:jtp-plus}
\sum_{n\in\IntegerNumbers}q^{\binom n2}z^n
=(q;q)_\infty(-z;q)_\infty(-q/z;q)_\infty.
\end{equation}
Equivalently,
\begin{equation}\label{eq:jtp-j}
 j(z;q):=(q,z,q/z;q)_\infty
 =\sum_{n\in\IntegerNumbers}(-1)^nq^{\binom n2}z^n.
\end{equation}
Both bilateral series converge normally on compact annuli in $\ComplexNumbers^\times$.
\end{theorem}
\begin{proof}
Put
\[
F(z)=(-z;q)_\infty(-q/z;q)_\infty.
\]
It is analytic on $\ComplexNumbers^\times$ and has a normally convergent Laurent expansion
$F(z)=\sum_{n\in\IntegerNumbers}c_nz^n$.  Directly from the products,
\[
F(qz)=\frac{1+z^{-1}}{1+z}F(z)=z^{-1}F(z).
\]
Comparing Laurent coefficients gives
$c_{n+1}=q^n c_n$, hence
$c_n=q^{\binom n2}c_0$ for every $n\in\IntegerNumbers$.

It remains to determine $c_0$.  Euler's expansion \eqref{eq:euler2} gives
\[
(-z;q)_\infty=\sum_{r\ge0}\frac{q^{\binom r2}z^r}{(q;q)_r},
\qquad
(-q/z;q)_\infty=\sum_{s\ge0}
 \frac{q^{s(s+1)/2}z^{-s}}{(q;q)_s}.
\]
Therefore
\[
c_0=\sum_{r=0}^{\infty}\frac{q^{r^2}}{(q;q)_r^2}.
\]
Take $c=q$ in the $q$-Gauss sum \eqref{eq:qgauss} and let
$a,b\to\infty$ after setting the argument equal to $q/(ab)$.  Termwise,
\[
\frac{(a,b;q)_r}{(q,q;q)_r}\left(\frac{q}{ab}\right)^r
\longrightarrow\frac{q^{r^2}}{(q;q)_r^2},
\]
and the product side tends to $1/(q;q)_\infty$.  Hence
$c_0=1/(q;q)_\infty$, which proves \eqref{eq:jtp-plus}.  Replacing $z$ by
$-z$ gives \eqref{eq:jtp-j}.
\end{proof}

\begin{corollary}[Euler pentagonal theorem]\label{cor:pentagonal}
For $|q|<1$,
\begin{equation}\label{eq:pentagonal}
(q;q)_\infty=\sum_{n\in\IntegerNumbers}(-1)^nq^{n(3n-1)/2}.
\end{equation}
\end{corollary}
\begin{proof}
In \eqref{eq:jtp-j}, replace the base by $q^3$ and set $z=q$.  The product is
$(q,q^2,q^3;q^3)_\infty=(q;q)_\infty$, and the exponent is
$3\binom n2+n=n(3n-1)/2$.
\end{proof}

\begin{theorem}[Jacobi quintuple product]\label{thm:qpp}
For $|q|<1$ and $z\ne0$,
\begin{align}
&\sum_{n\in\IntegerNumbers}q^{n(3n-1)/2}z^{3n}(1-zq^n)\label{eq:qpp}\\
&\qquad=(q,z,q/z;q)_\infty
(qz^2,q/z^2;q^2)_\infty.\notag
\end{align}
\end{theorem}
\begin{proof}
Let the product on the right be $P(z)$.  In terms of $j$ from
\eqref{eq:jtp-j},
\[
P(z)=\frac{j(z;q)j(qz^2;q^2)}{(q^2;q^2)_\infty}.
\]
The product gives the functional equation
\[
P(qz)=q^{-1}z^{-3}P(z).
\]
If $P(z)=\sum_{m\in\IntegerNumbers}c_mz^m$, then
\begin{equation}\label{eq:qpp-rec}
c_{m+3}=q^{m+1}c_m.
\end{equation}
Thus it is enough to determine $c_0,c_1,c_2$.

Using the two triple-product series,
\[
j(z;q)=\sum_{r\in\IntegerNumbers}(-1)^rq^{r(r-1)/2}z^r,
\qquad
j(qz^2;q^2)=\sum_{s\in\IntegerNumbers}(-1)^sq^{s^2}z^{2s}.
\]
Consequently
\begin{align*}
c_0&=\frac1{(q^2;q^2)_\infty}
 \sum_{s\in\IntegerNumbers}(-1)^sq^{3s^2+s}=1,\\
c_1&=-\frac1{(q^2;q^2)_\infty}
 \sum_{s\in\IntegerNumbers}(-1)^sq^{3s^2-s}=-1.
\end{align*}
Both evaluations are Euler's pentagonal theorem with base $q^2$, with
$s\mapsto-s$ in the first.  Finally,
\[
c_2=\frac1{(q^2;q^2)_\infty}
 \sum_{s\in\IntegerNumbers}(-1)^sq^{3s^2-3s+1}=0,
\]
because the terms indexed by $s$ and $1-s$ cancel.  The recurrence
\eqref{eq:qpp-rec} now gives
\[
c_{3n}=q^{n(3n-1)/2},\qquad
c_{3n+1}=-q^{n(3n+1)/2},\qquad c_{3n+2}=0,
\]
for all $n\in\IntegerNumbers$.  This is exactly the Laurent series on the left of
\eqref{eq:qpp}.
\end{proof}

\section{Partitions and coefficient recurrences}

\begin{definition}
A partition of $n$ is a weakly decreasing finite sequence of positive
integers with sum $n$.  Write $p(n)$ for the number of partitions of $n$, with
$p(0)=1$ and $p(n)=0$ for $n<0$.
\end{definition}

\begin{theorem}[Euler's partition generating function]\label{thm:partitiongf}
For $|q|<1$,
\begin{equation}\label{eq:partitiongf}
\sum_{n=0}^{\infty}p(n)q^n=\frac1{(q;q)_\infty}.
\end{equation}
\end{theorem}
\begin{proof}
For each possible part $m$, the geometric series
$1+q^m+q^{2m}+\cdots$ records its multiplicity.  Multiplying for all
$m\ge1$ gives
\[
\prod_{m\ge1}\frac1{1-q^m}.
\]
The coefficient of a fixed $q^n$ receives contributions from only finitely
many choices, so the argument is valid formally; for $|q|<1$ it is also
absolutely convergent.
\end{proof}

\begin{corollary}[Pentagonal recurrence]\label{cor:pentrec}
For $n\ge1$,
\begin{equation}\label{eq:pentrec}
p(n)=\sum_{k=1}^{\infty}(-1)^{k-1}
\left[p\!\left(n-\frac{k(3k-1)}2\right)
+p\!\left(n-\frac{k(3k+1)}2\right)\right],
\end{equation}
where the sum is finite because $p(m)=0$ for $m<0$.
\end{corollary}
\begin{proof}
Multiply \eqref{eq:partitiongf} by the pentagonal series
\eqref{eq:pentagonal}.  The product is $1$.  Equating the coefficient of
$q^n$ for $n>0$ and isolating $p(n)$ gives \eqref{eq:pentrec}.
\end{proof}

\begin{proposition}[Partitions in a rectangle]\label{prop:rectangle}
The coefficient of $q^m$ in $\GaussianBinomial{r+s}{r}{q}$ is the number of partitions of
$m$ having at most $r$ parts, each at most $s$.
\end{proposition}
\begin{proof}
Let $G_{r,s}(q)$ denote the indicated generating polynomial.  Separate the
partitions according to whether they have fewer than $r$ positive parts or
exactly $r$ positive parts.  In the second class subtract $1$ from every part;
this contributes a factor $q^r$ and leaves a partition with at most $r$ parts,
each at most $s-1$.  Hence
\[
G_{r,s}=G_{r-1,s}+q^rG_{r,s-1},
\]
with $G_{0,s}=G_{r,0}=1$.  By \eqref{eq:pascal2}, the Gaussian coefficient
$\GaussianBinomial{r+s}{r}{q}$ satisfies the same recurrence and boundary values.
\end{proof}

\begin{remark}
Many finite identities in the archive can be read in two complementary ways:
as terminating basic hypergeometric sums, or as identities between weighted
partition classes inside rectangles.  The analytic route is usually shorter;
the finite combinatorial route is often better for positivity.
\end{remark}

\section{Bailey pairs, inversion, and Bailey's lemma}

\begin{definition}[Bailey pair]
Two sequences $(\alpha_n)_{n\ge0}$ and $(\beta_n)_{n\ge0}$ form a
\emph{Bailey pair relative to $a$} when
\begin{equation}\label{eq:baileypair}
\beta_n=\sum_{r=0}^n
\frac{\alpha_r}{(q;q)_{n-r}(aq;q)_{n+r}}
\qquad(n\ge0).
\end{equation}
\end{definition}

\begin{lemma}[$q$-difference annihilation]\label{lem:qdiff}
If $P(x)$ is a polynomial of degree less than $N$, then
\begin{equation}\label{eq:qdiff}
\sum_{h=0}^N
\frac{(-1)^{N-h}q^{\binom{N-h}{2}}}{(q;q)_h(q;q)_{N-h}}
P(q^h)=0.
\end{equation}
\end{lemma}
\begin{proof}
By linearity it is enough to take $P(x)=x^m$, $0\le m<N$.  Put $r=N-h$.
The left-hand side becomes
\[
\frac{q^{mN}}{(q;q)_N}
\sum_{r=0}^N(-1)^rq^{\binom r2}\GaussianBinomial{N}{r}{q} q^{-mr}
=\frac{q^{mN}(q^{-m};q)_N}{(q;q)_N}
\]
by the finite $q$-binomial theorem.  Since $0\le m<N$, the product
$(q^{-m};q)_N$ contains the factor $1-q^0$ and is zero.
\end{proof}

\begin{theorem}[Bailey inversion]\label{thm:baileyinv}
The relation \eqref{eq:baileypair} is equivalent to
\begin{equation}\label{eq:baileyinv}
\alpha_n=\frac{1-aq^{2n}}{1-a}
\sum_{j=0}^n
\frac{(a;q)_{n+j}}{(q;q)_{n-j}}
(-1)^{n-j}q^{\binom{n-j}{2}}\beta_j.
\end{equation}
At $a=1$ the formula is interpreted by its removable limit.
\end{theorem}
\begin{proof}
Substitute \eqref{eq:baileypair} into the right-hand side of
\eqref{eq:baileyinv} and interchange the two finite sums.  The coefficient of
$\alpha_r$ is
\begin{equation}\label{eq:invkernel}
\frac{1-aq^{2n}}{1-a}
\sum_{j=r}^n
\frac{(a;q)_{n+j}(-1)^{n-j}q^{\binom{n-j}{2}}}
{(q;q)_{n-j}(q;q)_{j-r}(aq;q)_{j+r}}.
\end{equation}
If $r=n$, the sum has one term and its value is $1$.  Suppose $r<n$ and put
$N=n-r$, $h=j-r$.  The quotient of shifted factorials in
\eqref{eq:invkernel} is
\[
\frac{(a;q)_{n+j}}{(aq;q)_{j+r}}
=(1-a)(aq^{2r+h+1};q)_{N-1}.
\]
Apart from the harmless factor $1-aq^{2n}$, the coefficient is therefore the
left-hand side of \eqref{eq:qdiff} with
\[
P(x)=(aq^{2r+1}x;q)_{N-1},
\]
a polynomial of degree $N-1$.  It vanishes.  Thus the substitution recovers
$\alpha_n$ and no other $\alpha_r$, proving that the two triangular
relations are inverse.
\end{proof}

\begin{corollary}[The unit Bailey pair]\label{cor:unitpair}
Relative to $a$, the sequence $\beta_n=\delta_{n,0}$ has inverse
\begin{equation}\label{eq:unitalpha}
\alpha_n^{\mathrm{unit}}(a)
=\frac{1-aq^{2n}}{1-a}\frac{(a;q)_n}{(q;q)_n}
(-1)^nq^{\binom n2}.
\end{equation}
For $a=1$ this means
\begin{equation}\label{eq:unit-a1}
\alpha_0=1,\qquad
\alpha_n=(-1)^nq^{\binom n2}(1+q^n)\quad(n\ge1).
\end{equation}
\end{corollary}
\begin{proof}
Insert $\beta_0=1$ and $\beta_j=0$ for $j>0$ into
\eqref{eq:baileyinv}.  For $a\to1$, use
\[
\frac{(a;q)_n}{1-a}\longrightarrow(q;q)_{n-1},\qquad
\frac{1-q^{2n}}{1-q^n}=1+q^n.
\]
\end{proof}

\begin{theorem}[Bailey's lemma]\label{thm:baileylemma}
Suppose $(\alpha_n,\beta_n)$ is a Bailey pair relative to $a$.  Let
$\rho_1,\rho_2$ be auxiliary parameters and put
$A=aq/(\rho_1\rho_2)$.  Define
\begin{align}
\alpha_n'&=
\frac{(\rho_1,\rho_2;q)_n A^n}
{(aq/\rho_1,aq/\rho_2;q)_n}\alpha_n,
\label{eq:baileyalphaprime}\\
\beta_n'&=\sum_{j=0}^n
\frac{(\rho_1,\rho_2;q)_j(A;q)_{n-j}A^j}
{(q;q)_{n-j}(aq/\rho_1,aq/\rho_2;q)_n}\beta_j.
\label{eq:baileybetaprime}
\end{align}
Then $(\alpha_n',\beta_n')$ is again a Bailey pair relative to $a$.
\end{theorem}
\begin{proof}
Substitute \eqref{eq:baileypair} into \eqref{eq:baileybetaprime} and
interchange the finite $j$- and $r$-sums.  For a fixed $r$, put
$N=n-r$ and $j=r+s$.  Extract the factors independent of $s$ using
\eqref{eq:split}.  Reverse the factor $(A;q)_{N-s}$ by
\eqref{eq:reverse}.  The remaining sum over $s$ is
\[
{}_3\phi_2\!\left(
\begin{matrix}\rho_1q^r,\rho_2q^r,q^{-N}\\
aq^{2r+1},\rho_1\rho_2q^{-N}/a
\end{matrix};q,q\right).
\]
It is balanced because the second lower parameter is
\[
\frac{(\rho_1q^r)(\rho_2q^r)q^{1-N}}{aq^{2r+1}}
=\frac{\rho_1\rho_2q^{-N}}a.
\]
The $q$-Pfaff--Saalsch\"utz sum \eqref{eq:qps} evaluates it.  Substituting
that product and cancelling the extracted shifted factorials leaves
\[
\frac1{(q;q)_{n-r}(aq;q)_{n+r}}
\frac{(\rho_1,\rho_2;q)_rA^r}
{(aq/\rho_1,aq/\rho_2;q)_r}\alpha_r.
\]
Summing over $r$ gives exactly
\[
\beta_n'=\sum_{r=0}^n
\frac{\alpha_r'}{(q;q)_{n-r}(aq;q)_{n+r}},
\]
which is the required Bailey-pair relation.
\end{proof}

\begin{corollary}[Bailey-chain step]\label{cor:baileychain}
If $(\alpha,\beta)$ is a Bailey pair relative to $a$, then so is
\begin{align}
\widehat\alpha_n&=a^nq^{n^2}\alpha_n,\label{eq:chainalpha}\\
\widehat\beta_n&=\sum_{j=0}^n
\frac{a^jq^{j^2}}{(q;q)_{n-j}}\beta_j.
\label{eq:chainbeta}
\end{align}
\end{corollary}
\begin{proof}
Let $\rho_1,\rho_2\to\infty$ in Bailey's lemma.  For fixed $n$,
\[
(\rho_1,\rho_2;q)_n
\left(\frac{aq}{\rho_1\rho_2}\right)^n
\longrightarrow a^nq^{n^2},
\]
while all shifted factorials whose arguments tend to zero tend to $1$.
The formulas reduce to \eqref{eq:chainalpha}--\eqref{eq:chainbeta}.
\end{proof}

\begin{corollary}[Weak Bailey lemma]\label{cor:weakbailey}
Under absolute convergence,
\begin{equation}\label{eq:weakbailey}
\sum_{n=0}^{\infty}a^nq^{n^2}\beta_n
=\frac1{(aq;q)_\infty}
\sum_{n=0}^{\infty}a^nq^{n^2}\alpha_n.
\end{equation}
\end{corollary}
\begin{proof}
Apply the chain step and write the defining Bailey relation for
$(\widehat\alpha,\widehat\beta)$ at index $N$.  Let $N\to\infty$.  From
\eqref{eq:chainbeta},
\[
\widehat\beta_N\longrightarrow
\frac1{(q;q)_\infty}\sum_{j\ge0}a^jq^{j^2}\beta_j.
\]
From the Bailey-pair relation,
\[
\widehat\beta_N\longrightarrow
\frac1{(q;q)_\infty(aq;q)_\infty}
\sum_{r\ge0}a^rq^{r^2}\alpha_r.
\]
Dominated convergence is justified whenever the two displayed series are
absolutely convergent.  Multiplication by $(q;q)_\infty$ gives
\eqref{eq:weakbailey}.
\end{proof}

\begin{lemma}[Parameter lowering from $a=q$ to $a=1$]\label{lem:lowering}
Suppose $(\alpha_n,\beta_n)$ is a Bailey pair relative to $q$.  Define
\begin{align}
\widetilde\alpha_0&=\alpha_0,\label{eq:lower0}\\
\widetilde\alpha_n&=(1-q)
\left(\frac{\alpha_n}{1-q^{2n+1}}
-\frac{q^{2n-1}\alpha_{n-1}}{1-q^{2n-1}}\right),
\qquad n\ge1.\label{eq:lowern}
\end{align}
Then $(\widetilde\alpha_n,\beta_n)$ is a Bailey pair relative to $1$.
\end{lemma}
\begin{proof}
Let
\[
K_{n,r}=\frac1{(q;q)_{n-r}(q;q)_{n+r}}.
\]
In
$\sum_{r=0}^n\widetilde\alpha_rK_{n,r}$, the coefficient of a fixed
$\alpha_r$ is
\[
\frac{1-q}{1-q^{2r+1}}K_{n,r}
-\frac{(1-q)q^{2r+1}}{1-q^{2r+1}}K_{n,r+1},
\]
where the second term is absent for $r=n$.  Since
\[
\frac{K_{n,r+1}}{K_{n,r}}
=\frac{1-q^{n-r}}{1-q^{n+r+1}},
\]
the coefficient simplifies to
\[
\frac{1-q}{1-q^{n+r+1}}K_{n,r}
=\frac1{(q;q)_{n-r}(q^2;q)_{n+r}}.
\]
This is exactly the kernel of a Bailey pair relative to $a=q$.  Therefore
\[
\sum_{r=0}^n\frac{\widetilde\alpha_r}
{(q;q)_{n-r}(q;q)_{n+r}}
=\sum_{r=0}^n\frac{\alpha_r}
{(q;q)_{n-r}(q^2;q)_{n+r}}
=\beta_n.
\]
\end{proof}

\begin{remark}
The elementary lowering step is the missing bridge between the two endpoint
Bailey chains $a=1$ and $a=q$.  Placing it at different depths in a chain
produces the linear terms in all members of the Andrews--Gordon family.
\end{remark}

\section{The Rogers--Ramanujan identities}

\begin{theorem}[Rogers--Ramanujan identities]\label{thm:RR}
For $|q|<1$,
\begin{align}
\sum_{n=0}^{\infty}\frac{q^{n^2}}{(q;q)_n}
&=\frac1{(q,q^4;q^5)_\infty},\label{eq:RR1}\\
\sum_{n=0}^{\infty}\frac{q^{n^2+n}}{(q;q)_n}
&=\frac1{(q^2,q^3;q^5)_\infty}.\label{eq:RR2}
\end{align}
\end{theorem}
\begin{proof}
Start with the unit Bailey pair \cref{cor:unitpair}.  Apply one Bailey-chain
step.  Because the original $\beta$ is $\delta_{n,0}$, the new pair has
\[
\widehat\beta_n=\frac1{(q;q)_n},\qquad
\widehat\alpha_n=a^nq^{n^2}\alpha_n^{\mathrm{unit}}(a).
\]
Apply the weak Bailey lemma to this pair:
\begin{equation}\label{eq:RRmaster}
\sum_{n\ge0}\frac{a^nq^{n^2}}{(q;q)_n}
=\frac1{(aq;q)_\infty}
\sum_{n\ge0}a^{2n}q^{2n^2}\alpha_n^{\mathrm{unit}}(a).
\end{equation}

First let $a\to1$.  By \eqref{eq:unit-a1}, the numerator series on the right
is
\begin{align*}
1+\sum_{n\ge1}(-1)^n
\left(q^{(5n^2-n)/2}+q^{(5n^2+n)/2}\right)
&=\sum_{m\in\IntegerNumbers}(-1)^mq^{m(5m-1)/2}\\
&=(q^2,q^3,q^5;q^5)_\infty,
\end{align*}
where the last equality is Jacobi's triple product with base $q^5$ and
$z=q^2$.  Division by
\[
(q;q)_\infty=(q,q^2,q^3,q^4,q^5;q^5)_\infty
\]
gives \eqref{eq:RR1}.

Now put $a=q$.  Formula \eqref{eq:unitalpha} becomes
\[
\alpha_n^{\mathrm{unit}}(q)
=(-1)^nq^{\binom n2}\frac{1-q^{2n+1}}{1-q}.
\]
Since $(q^2;q)_\infty=(q;q)_\infty/(1-q)$, the right-hand side of
\eqref{eq:RRmaster} is
\[
\frac1{(q;q)_\infty}
\sum_{m\in\IntegerNumbers}(-1)^mq^{m(5m+3)/2}.
\]
Indeed, expanding $1-q^{2n+1}$ supplies the nonnegative and negative halves
of the bilateral sum, with the second term corresponding to $m=-n-1$.
The triple product with base $q^5$ and $z=q^4$ evaluates the bilateral sum as
$(q,q^4,q^5;q^5)_\infty$.  Cancelling residue classes in $(q;q)_\infty$
gives \eqref{eq:RR2}.
\end{proof}

\begin{remark}[Two meanings of the products]
The product in \eqref{eq:RR1} generates partitions into parts congruent to
$1$ or $4$ modulo $5$; the product in \eqref{eq:RR2} generates partitions
into parts congruent to $2$ or $3$ modulo $5$.  The classical difference-two
partition interpretations require an additional combinatorial argument.  The
analytic identities themselves are completely proved above.
\end{remark}

\section{The Rogers--Ramanujan continued fraction}

For $m\ge0$ define
\begin{equation}\label{eq:Fm}
F_m(q)=\sum_{n=0}^{\infty}\frac{q^{n^2+mn}}{(q;q)_n}.
\end{equation}

\begin{lemma}[Continuant recurrence]\label{lem:Frec}
For $m\ge0$ and $|q|<1$,
\begin{equation}\label{eq:Frec}
F_m=F_{m+1}+q^{m+1}F_{m+2}.
\end{equation}
\end{lemma}
\begin{proof}
Subtract the two series $F_m-F_{m+1}$.  The $n=0$ terms cancel, and for
$n\ge1$ use $(1-q^n)/(q;q)_n=1/(q;q)_{n-1}$:
\begin{align*}
F_m-F_{m+1}
&=\sum_{n\ge1}\frac{q^{n^2+mn}(1-q^n)}{(q;q)_n}\\
&=\sum_{r\ge0}\frac{q^{(r+1)^2+m(r+1)}}{(q;q)_r}
=q^{m+1}F_{m+2}.
\end{align*}
\end{proof}

\begin{theorem}[Rogers--Ramanujan continued fraction]\label{thm:RRCF}
For $|q|<1$,
\begin{align}
\frac{F_1(q)}{F_0(q)}
&=\cfrac1{1+\cfrac{q}{1+\cfrac{q^2}{1+\cfrac{q^3}{1+\ddots}}}},
\label{eq:RRcf-series}\\
R(q):=q^{1/5}\frac{F_1(q)}{F_0(q)}
&=q^{1/5}
\frac{(q,q^4;q^5)_\infty}{(q^2,q^3;q^5)_\infty}.
\label{eq:RRcf-product}
\end{align}
\end{theorem}
\begin{proof}
Put $C_m=F_{m+1}/F_m$.  The recurrence \eqref{eq:Frec} gives
\begin{equation}\label{eq:Crec}
C_m=\frac1{1+q^{m+1}C_{m+1}}.
\end{equation}
Iterating to depth $M$ gives a finite continuant with tail $C_M$.  Moreover,
$F_m\to1$ as $m\to\infty$: indeed, for $n\ge1$,
\[
\left|\frac{q^{n^2+mn}}{(q;q)_n}\right|
\le\frac{|q|^{n^2+mn}}{(|q|;|q|)_\infty},
\]
and the majorant tends summably to zero.  Hence $C_M\to1$.  Taking the limit
in the finite continuants proves \eqref{eq:RRcf-series}, including
convergence.  Finally \cref{thm:RR} gives
\[
\frac{F_1}{F_0}
=\frac{(q,q^4;q^5)_\infty}{(q^2,q^3;q^5)_\infty},
\]
which is \eqref{eq:RRcf-product}.
\end{proof}

\section{The full Andrews--Gordon family}

\begin{theorem}[Andrews--Gordon identities]\label{thm:AG}
Let $k\ge2$ and $1\le i\le k$.  For
$n_1,\ldots,n_{k-1}\ge0$ put
\[
N_j=n_j+n_{j+1}+\cdots+n_{k-1}\qquad(1\le j\le k-1).
\]
Then, for $|q|<1$,
\begin{align}
&\sum_{n_1,\ldots,n_{k-1}\ge0}
\frac{q^{N_1^2+\cdots+N_{k-1}^2+N_i+\cdots+N_{k-1}}}
{(q;q)_{n_1}\cdots(q;q)_{n_{k-1}}}
\label{eq:AG}\\
&\qquad=
\frac{(q^i,q^{2k+1-i},q^{2k+1};q^{2k+1})_\infty}
{(q;q)_\infty}.
\notag
\end{align}
When $i=k$, the linear sum $N_i+\cdots+N_{k-1}$ is empty and equals zero.
\end{theorem}
\begin{proof}
We keep every Bailey operation visible.  A chain step applied repeatedly to a
pair shows by induction that nested indices are introduced in weakly
decreasing order.  After $r$ chain steps and a final weak limit, the
$\beta$-side has $r$ indices
$M_1\ge\cdots\ge M_r\ge0$, denominator
\[
(q;q)_{M_1-M_2}\cdots(q;q)_{M_{r-1}-M_r}(q;q)_{M_r},
\]
and one factor $a^{M_j}q^{M_j^2}$ for every operation carried out while the
relative parameter is $a$.  This statement follows immediately from
\eqref{eq:chainbeta}; it is also proved by induction on $r$.

\smallskip
\noindent\emph{Case $i=1$.}
Start with the unit pair relative to $q$, perform $k-1$ chain steps relative
to $q$, and apply the weak Bailey lemma relative to $q$.  All $k-1$ nested
indices acquire a linear factor $q^{M_j}$.  With
$n_j=M_j-M_{j+1}$ and $M_k=0$, the $\beta$-side is the left-hand side of
\eqref{eq:AG} for $i=1$.

On the $\alpha$-side, the $k-1$ chain steps and the weak limit multiply the
unit alpha by $q^{k(n^2+n)}$.  Since
$(q^2;q)_\infty=(q;q)_\infty/(1-q)$, the result is
\begin{align*}
\frac1{(q;q)_\infty}
\sum_{m\in\IntegerNumbers}(-1)^m
q^{((2k+1)m^2+(2k-1)m)/2}.
\end{align*}
The negative half is obtained from the term containing
$-q^{2n+1}$ by the substitution $m=-n-1$.  Jacobi's triple product, with
base $q^{2k+1}$ and $z=q^{2k}$, evaluates the bilateral sum as
\[
(q,q^{2k},q^{2k+1};q^{2k+1})_\infty,
\]
which proves the required product for $i=1$.

\smallskip
\noindent\emph{Cases $2\le i\le k$.}
Begin again with the unit pair relative to $q$.  Perform
\[
c=k-i+1
\]
chain steps relative to $q$, apply the lowering lemma
\cref{lem:lowering}, then perform $i-2$ chain steps relative to $1$, and
finally apply the weak Bailey lemma relative to $1$.  There are
$c+(i-2)=k-1$ chain steps in total.  The lowering operation does not change
$\beta$.  Exactly the deepest $c-1=k-i$ indices were weighted by the
$q$-parameter before lowering; after the outer $1$-steps and weak limit they
are $M_i,\ldots,M_{k-1}$.  Thus the $\beta$-side is precisely the left-hand
side of \eqref{eq:AG}.

It remains to compute the transformed alpha sequence.  Before lowering it is
\[
\alpha_n=(-1)^nq^{\binom n2+c(n^2+n)}
\frac{1-q^{2n+1}}{1-q}.
\]
For $n\ge1$, \eqref{eq:lowern} gives
\begin{align*}
\widetilde\alpha_n=(-1)^n\biggl(&
q^{n(n-1)/2+c(n^2+n)}
+q^{n(n+1)/2+c(n^2-n)}\biggr),
\end{align*}
while $\widetilde\alpha_0=1$.  The remaining $i-2$ chain steps and the weak
limit multiply this by $q^{(i-1)n^2}$.  Since
$c+i-1=k$, the resulting one-sided sum is the two halves of
\begin{equation}\label{eq:AGbilateral}
\sum_{m\in\IntegerNumbers}(-1)^m
q^{((2k+1)m^2+(2k-2i+1)m)/2}.
\end{equation}
Therefore the multiple sum equals \eqref{eq:AGbilateral} divided by
$(q;q)_\infty$.  Apply the triple product with base $q^{2k+1}$ and
$z=q^{2k+1-i}$.  Its product is
\[
(q^{2k+1},q^{2k+1-i},q^i;q^{2k+1})_\infty,
\]
which proves \eqref{eq:AG}.
\end{proof}

\begin{corollary}[Rogers--Ramanujan as modulus five]\label{cor:AG-RR}
Taking $k=2$ and $i=2$ or $i=1$ in \cref{thm:AG} gives
\eqref{eq:RR1} and \eqref{eq:RR2}, respectively.
\end{corollary}
\begin{proof}
There is one variable, $N_1=n_1$.  For $i=2$ the linear term is absent; for
$i=1$ it is $n_1$.  The products reduce to the two residue-class products in
\cref{thm:RR}.
\end{proof}

\begin{corollary}[The modulus seven family]\label{cor:mod7}
For $r,s\ge0$, set $N_1=r+s$, $N_2=s$.  Then the three sums with exponents
\[
N_1^2+N_2^2+N_1+N_2,\qquad
N_1^2+N_2^2+N_2,\qquad
N_1^2+N_2^2
\]
are respectively the products obtained by excluding residues
$0,\pm1$, $0,\pm2$, and $0,\pm3$ modulo $7$.
\end{corollary}
\begin{proof}
This is \cref{thm:AG} with $k=3$ and $i=1,2,3$, after cancelling the three
listed residue classes from $(q;q)_\infty$.
\end{proof}

\begin{remark}[Moduli $9$ and $27$]
The same theorem with $k=4$ yields all four odd-modulus Gordon products of
modulus $9$, and $k=13$ yields the thirteen products of modulus $27$.  This
places a large part of the modulus-specific archive under a single proof.
Even modulus $14$ belongs to the Bressoud-type extension and requires a
parity-sensitive Bailey step; it is discussed in \cref{sec:extensions}.
\end{remark}

\section{Poisson summation and theta transformations}

\begin{definition}[Fourier transform]
For a Schwartz function $f\colon\RealNumbers\to\ComplexNumbers$, use the normalization
\[
\widehat f(\xi)=\int_{-\infty}^{\infty}f(x)\EulerE^{-2\pi\ImaginaryUnit x\xi}\,dx.
\]
\end{definition}

\begin{theorem}[Poisson summation]\label{thm:poisson}
For every Schwartz function $f$,
\begin{equation}\label{eq:poisson}
\sum_{n\in\IntegerNumbers}f(n)=\sum_{m\in\IntegerNumbers}\widehat f(m).
\end{equation}
Both sums converge absolutely.
\end{theorem}
\begin{proof}
Periodize $f$:
\[
P_f(x)=\sum_{n\in\IntegerNumbers}f(x+n).
\]
Schwartz decay makes this series, and every termwise derivative, uniformly
convergent on compact sets.  Thus $P_f$ is a smooth $1$-periodic function.
Its $m$th Fourier coefficient is, after an absolutely convergent interchange,
\[
\int_0^1P_f(x)\EulerE^{-2\pi\ImaginaryUnit mx}\,dx
=\int_{\RealNumbers}f(x)\EulerE^{-2\pi\ImaginaryUnit mx}\,dx=\widehat f(m).
\]
Repeated integration by parts shows that these coefficients decay faster
than any power of $|m|$, so the Fourier series converges absolutely and
uniformly to $P_f$.  Evaluating it at $x=0$ gives \eqref{eq:poisson}.
\end{proof}

\begin{lemma}[Fourier transform of a Gaussian]\label{lem:gaussianFT}
For $t>0$,
\begin{equation}\label{eq:gaussianFT}
\int_{\RealNumbers}\EulerE^{-\pi t x^2}\EulerE^{-2\pi\ImaginaryUnit x\xi}\,dx
=t^{-1/2}\EulerE^{-\pi\xi^2/t}.
\end{equation}
\end{lemma}
\begin{proof}
For $t=1$, let the integral be $G(\xi)$.  Differentiating under the integral
and integrating by parts gives
$G'(\xi)=-2\pi\xi G(\xi)$.  Hence
$G(\xi)=G(0)\EulerE^{-\pi\xi^2}$.  The classical polar-coordinate computation
\[
G(0)^2=\iint_{\RealNumbers^2}\EulerE^{-\pi(x^2+y^2)}\,dx\,dy
=2\pi\int_0^\infty r\EulerE^{-\pi r^2}\,dr=1
\]
shows $G(0)=1$.  Scaling $x\mapsto x/\sqrt t$ proves
\eqref{eq:gaussianFT}.
\end{proof}

For $\tau$ in the upper half-plane define
\begin{align*}
\vartheta_2(\tau)&=\sum_{n\in\IntegerNumbers}\EulerE^{\pi\ImaginaryUnit\tau(n+1/2)^2},\\
\vartheta_3(\tau)&=\sum_{n\in\IntegerNumbers}\EulerE^{\pi\ImaginaryUnit\tau n^2},\\
\vartheta_4(\tau)&=\sum_{n\in\IntegerNumbers}(-1)^n\EulerE^{\pi\ImaginaryUnit\tau n^2}.
\end{align*}
We take $(-\ImaginaryUnit\tau)^{1/2}$ to be the holomorphic square root that is positive
for $\tau=\ImaginaryUnit t$, $t>0$.

\begin{theorem}[Theta inversion]\label{thm:thetaS}
For $\operatorname{Im}\tau>0$,
\begin{align}
\vartheta_3(-1/\tau)&=(-\ImaginaryUnit\tau)^{1/2}\vartheta_3(\tau),\label{eq:theta3S}\\
\vartheta_2(-1/\tau)&=(-\ImaginaryUnit\tau)^{1/2}\vartheta_4(\tau),\label{eq:theta2S}\\
\vartheta_4(-1/\tau)&=(-\ImaginaryUnit\tau)^{1/2}\vartheta_2(\tau).\label{eq:theta4S}
\end{align}
\end{theorem}
\begin{proof}
First set $\tau=\ImaginaryUnit t$, $t>0$.  Apply Poisson summation to
$f(x)=\EulerE^{-\pi t x^2}$ and use \eqref{eq:gaussianFT}:
\[
\sum_{n\in\IntegerNumbers}\EulerE^{-\pi t n^2}
=t^{-1/2}\sum_{m\in\IntegerNumbers}\EulerE^{-\pi m^2/t}.
\]
This is \eqref{eq:theta3S} on the positive imaginary axis.  Applying Poisson
summation to $f(x+1/2)$ introduces the Fourier factor $(-1)^m$ and gives
\eqref{eq:theta2S}; applying it to $(-1)^nf(n)$ is the inverse relation
\eqref{eq:theta4S}.  Each side is holomorphic in the upper half-plane and the
chosen square root is holomorphic there.  The identity theorem extends all
three equalities from the positive imaginary axis to the whole upper
half-plane.
\end{proof}

\begin{proposition}[Theta products]\label{prop:thetaproducts}
Put $Q=\EulerE^{\pi\ImaginaryUnit\tau}$.  Then
\begin{align}
\vartheta_2(\tau)&=2Q^{1/4}
\prod_{n=1}^{\infty}(1-Q^{2n})(1+Q^{2n})^2,\label{eq:theta2prod}\\
\vartheta_3(\tau)&=
\prod_{n=1}^{\infty}(1-Q^{2n})(1+Q^{2n-1})^2,\label{eq:theta3prod}\\
\vartheta_4(\tau)&=
\prod_{n=1}^{\infty}(1-Q^{2n})(1-Q^{2n-1})^2.\label{eq:theta4prod}
\end{align}
\end{proposition}
\begin{proof}
These are direct specializations of Jacobi's triple product.  For example,
use base $Q^2$ in \eqref{eq:jtp-plus}; setting $z=1$ gives
\eqref{eq:theta3prod}, and setting $z=-1$ gives
\eqref{eq:theta4prod}.  For $\vartheta_2$, factor
$Q^{(n+1/2)^2}=Q^{1/4}Q^{n(n+1)}$ and pair the terms indexed by $n$ and
$-n-1$ before applying the same product with a shifted argument.  The result
is \eqref{eq:theta2prod}.
\end{proof}

\section{The Dedekind eta function and the $q\to1^-$ regime}

\begin{definition}[Dedekind eta]
For $\operatorname{Im}\tau>0$,
\[
\eta(\tau)=\EulerE^{\pi\ImaginaryUnit\tau/12}
\prod_{n=1}^{\infty}(1-\EulerE^{2\pi\ImaginaryUnit n\tau}).
\]
\end{definition}

\begin{lemma}[Eta and the three theta constants]\label{lem:eta-theta}
For $\operatorname{Im}\tau>0$,
\begin{equation}\label{eq:eta-theta}
2\eta(\tau)^3=\vartheta_2(\tau)\vartheta_3(\tau)\vartheta_4(\tau).
\end{equation}
\end{lemma}
\begin{proof}
Multiply the three product formulas in \cref{prop:thetaproducts}.  With
$q=Q^2=\EulerE^{2\pi\ImaginaryUnit\tau}$, the extra factors besides
$2q^{1/8}(q;q)_\infty^3$ cancel because
\[
(-q;q)_\infty(q;q^2)_\infty
=\frac{(q^2;q^2)_\infty}{(q;q)_\infty}(q;q^2)_\infty=1.
\]
Since $\eta(\tau)^3=q^{1/8}(q;q)_\infty^3$, the result follows.
\end{proof}

\begin{theorem}[Eta transformations]\label{thm:eta-transform}
For $\operatorname{Im}\tau>0$,
\begin{align}
\eta(\tau+1)&=\EulerE^{\pi\ImaginaryUnit/12}\eta(\tau),\label{eq:etaT}\\
\eta(-1/\tau)&=(-\ImaginaryUnit\tau)^{1/2}\eta(\tau).
\label{eq:etaS}
\end{align}
\end{theorem}
\begin{proof}
Equation \eqref{eq:etaT} follows directly from the definition, since the
infinite product is unchanged and the leading exponential is multiplied by
$\EulerE^{\pi\ImaginaryUnit/12}$.

For inversion, combine \eqref{eq:eta-theta} with the three theta
transformations.  Their product gives
\[
\eta(-1/\tau)^3=(-\ImaginaryUnit\tau)^{3/2}\eta(\tau)^3.
\]
Thus the quotient of the two sides of \eqref{eq:etaS} is a holomorphic
cube root of unity, hence constant.  On the positive imaginary axis all
products are positive real and the chosen square root is positive, so the
constant is $1$.
\end{proof}

\begin{theorem}[Exact modular form of $(q;q)_\infty$]\label{thm:qpoch-asymp}
Let $t>0$ and $q=\EulerE^{-t}$.  Then
\begin{equation}\label{eq:qpoch-exact}
(q;q)_\infty
=\sqrt{\frac{2\pi}{t}}
\exp\!\left(-\frac{\pi^2}{6t}+\frac{t}{24}\right)
\bigl(\EulerE^{-4\pi^2/t};\EulerE^{-4\pi^2/t}\bigr)_\infty.
\end{equation}
Consequently, as $t\to0^+$,
\begin{align}
(q;q)_\infty
&=\sqrt{\frac{2\pi}{t}}
\exp\!\left(-\frac{\pi^2}{6t}+\frac{t}{24}\right)
\left(1+O(\EulerE^{-4\pi^2/t})\right),\label{eq:qpoch-asymp}\\
\log(q;q)_\infty
&=-\frac{\pi^2}{6t}+\frac12\log\frac{2\pi}{t}
+\frac{t}{24}+O(\EulerE^{-4\pi^2/t}).\label{eq:logqpoch}
\end{align}
\end{theorem}
\begin{proof}
Put $\tau=\ImaginaryUnit t/(2\pi)$.  From the definition,
\[
\eta(\tau)=\EulerE^{-t/24}(\EulerE^{-t};\EulerE^{-t})_\infty.
\]
The eta inversion formula gives
\[
\eta\!\left(\frac{\ImaginaryUnit t}{2\pi}\right)
=\sqrt{\frac{2\pi}{t}}\,
\eta\!\left(\frac{2\pi\ImaginaryUnit}{t}\right).
\]
Expanding the eta function on the right yields
\[
\eta\!\left(\frac{2\pi\ImaginaryUnit}{t}\right)
=\EulerE^{-\pi^2/(6t)}
(\EulerE^{-4\pi^2/t};\EulerE^{-4\pi^2/t})_\infty.
\]
Combining these identities proves \eqref{eq:qpoch-exact}.  The remaining
product differs from $1$ by $O(\EulerE^{-4\pi^2/t})$; taking a logarithm preserves
that error order and gives \eqref{eq:qpoch-asymp}--\eqref{eq:logqpoch}.
\end{proof}

\begin{remark}
Formula \eqref{eq:qpoch-exact}, not merely its leading asymptotic term, is a
useful numerical acceleration: when $q$ is close to $1$, the transformed nome
$\exp(-4\pi^2/t)$ is exponentially small.
\end{remark}

\section{Ramanujan's bilateral ${}_1\psi_1$ summation}

For an integer $n<0$, extend the shifted factorial by
\[
(a;q)_n=\frac1{(aq^n;q)_{-n}}.
\]
This agrees with $(a;q)_n=(a;q)_\infty/(aq^n;q)_\infty$ whenever both sides
are defined.

\begin{theorem}[Ramanujan's ${}_1\psi_1$ sum]\label{thm:1psi1}
If $|q|<1$ and $|b/a|<|z|<1$, then
\begin{equation}\label{eq:1psi1}
\sum_{n\in\IntegerNumbers}\frac{(a;q)_n}{(b;q)_n}z^n
=\frac{(q,b/a,az,q/(az);q)_\infty}
{(b,q/a,z,b/(az);q)_\infty}.
\end{equation}
\end{theorem}
\begin{proof}
We first prove the formula for $b=q^{m+1}$, $m\in\NaturalNumbers$.  For these values,
terms with $n<-m$ vanish, because the denominator has a pole.  Put $n=k-m$.
Then
\begin{align*}
\sum_{n\in\IntegerNumbers}\frac{(a;q)_n}{(q^{m+1};q)_n}z^n
&=z^{-m}\frac{(a;q)_{-m}}{(q^{m+1};q)_{-m}}
\sum_{k=0}^{\infty}
\frac{(aq^{-m};q)_k}{(q;q)_k}z^k\\
&=z^{-m}\frac{(q;q)_m}{(aq^{-m};q)_m}
\frac{(aq^{-m}z;q)_\infty}{(z;q)_\infty},
\end{align*}
where the last step is the $q$-binomial theorem.  Repeated use of
\eqref{eq:reverse} and \eqref{eq:split} simplifies this expression to
\[
\frac{(q,q^{m+1}/a,az,q/(az);q)_\infty}
{(q^{m+1},q/a,z,q^{m+1}/(az);q)_\infty},
\]
which is the right-hand side of \eqref{eq:1psi1} at $b=q^{m+1}$.

Fix $a,z,q$ with $0<|z|<1$.  On every compact subset of
$|b|<|az|$, the bilateral series converges normally: the positive tail is
controlled by $|z|^n$, while after writing $n=-r$ and using
\eqref{eq:reverse}, the negative tail is controlled by
$|b/(az)|^r$.  Hence both sides of \eqref{eq:1psi1} are holomorphic functions
of $b$ in that disk, away from removable product singularities.  The values
$b=q^{m+1}$ accumulate at $0$, so the identity theorem extends the equality
to $|b|<|az|$, which is exactly $|b/a|<|z|$.
\end{proof}

\begin{corollary}[Bilateral-to-product principle]\label{cor:bilateralprinciple}
Any specialization of \eqref{eq:1psi1} obtained by a limit that is locally
uniform on a compact subannulus may be taken termwise.  In particular,
Jacobi-type bilateral theta sums arise by sending one parameter to $0$ or
$\infty$ after a compensating rescaling of $z$.
\end{corollary}
\begin{proof}
Normal convergence on compact subannuli permits termwise limits by the
Weierstrass theorem.  The product side is handled by local uniform convergence
of the infinite products.
\end{proof}

\section{Finite forms and polynomial approximants}

Infinite product identities are often best understood through finite
polynomials.  Besides giving rigorous convergence, finite forms retain
coefficient positivity and are suitable for computer-algebra certification.

\begin{proposition}[A finite Rogers--Ramanujan recurrence]\label{prop:finiteRR}
Let $A_m(q),B_m(q)$ be the numerator and denominator continuants of
\[
\cfrac1{1+\cfrac{q}{1+\cfrac{q^2}{\ddots+\cfrac{q^m}{1}}}}.
\]
With initial data $A_{-1}=1,A_0=1$ and $B_{-1}=0,B_0=1$, they satisfy
\begin{equation}\label{eq:continuants}
A_m=A_{m-1}+q^mA_{m-2},\qquad
B_m=B_{m-1}+q^mB_{m-2}.
\end{equation}
The finite fraction is $A_{m-1}/A_m$ after a consistent choice of indexing,
and it converges to $F_1/F_0$ as $m\to\infty$.
\end{proposition}
\begin{proof}
The continuant recurrence follows by clearing the last denominator in a
finite continued fraction.  Induction proves the stated numerator-denominator
recurrences.  The tail argument in \cref{thm:RRCF} shows that replacing the
true tail $C_m$ by $1$ changes the value by a quantity tending to zero, so the
finite continuants converge to $F_1/F_0$.
\end{proof}

\begin{proposition}[Coefficientwise passage to infinite products]\label{prop:coefflimit}
Let $P_N(q)$ be a sequence of polynomials such that, for every fixed degree
$d$, the coefficients of $q^0,\ldots,q^d$ stabilize as $N\to\infty$.  Then
$P_N$ has a well-defined coefficientwise limit in $\ComplexNumbers[[q]]$.  If in addition
$P_N\to P$ locally uniformly on $|q|<1$, the formal limit is the Taylor series
of $P$.
\end{proposition}
\begin{proof}
Coefficient stabilization defines a unique formal series.  Local uniform
convergence allows coefficient extraction on a small circle by Cauchy's
integral formula; the coefficient limits must equal the Taylor coefficients
of $P$.
\end{proof}

\begin{remark}
This elementary principle is the rigorous link between finite Gaussian
polynomial identities and their infinite Rogers--Ramanujan or
Andrews--Gordon limits.  It also prevents a common error in experimental
work: pointwise numerical convergence alone does not imply coefficientwise
stabilization.
\end{remark}

\section{Cyclotomic dissections and Borwein-type sign questions}
\label{sec:borwein}

Define the finite product
\begin{equation}\label{eq:borweinP}
P_n(q)=\frac{(q;q)_{3n}}{(q^3;q^3)_n}
=\prod_{j=1}^n(1-q^{3j-2})(1-q^{3j-1}).
\end{equation}
There are unique polynomials $A_n,B_n,C_n\in\IntegerNumbers[x]$ such that
\begin{equation}\label{eq:borweindissection}
P_n(q)=A_n(q^3)-qB_n(q^3)-q^2C_n(q^3).
\end{equation}
Uniqueness follows by collecting exponents modulo $3$.

\begin{proposition}[Reciprocity constraints]\label{prop:borweinreciprocity}
The three residue polynomials satisfy
\begin{align}
A_n(x)&=x^{n^2}A_n(x^{-1}),\label{eq:Arecip}\\
B_n(x)&=x^{n^2-1}C_n(x^{-1}),\label{eq:Brecip}\\
C_n(x)&=x^{n^2-1}B_n(x^{-1}).\label{eq:Crecip}
\end{align}
In particular, $B_n(1)=C_n(1)$ and $A_n(1)=2B_n(1)$ for $n\ge1$.
\end{proposition}
\begin{proof}
The degree of $P_n$ is
\[
\sum_{j=1}^n[(3j-2)+(3j-1)]=3n^2.
\]
Each pair of factors is reciprocal with positive sign, so
$P_n(q)=q^{3n^2}P_n(q^{-1})$.  Substitute
\eqref{eq:borweindissection} on both sides and write $x=q^3$.  Comparing the
three residue classes modulo $3$ gives \eqref{eq:Arecip}--\eqref{eq:Crecip}.
Putting $x=1$ gives $B_n(1)=C_n(1)$.  Finally $P_n(1)=0$ and
\eqref{eq:borweindissection} yield $A_n(1)=B_n(1)+C_n(1)$.
\end{proof}

\begin{conjecture}[Borwein nonnegativity]\label{conj:borwein}
For every $n\ge1$, all coefficients of $A_n(x)$, $B_n(x)$, and $C_n(x)$ in
\eqref{eq:borweindissection} are nonnegative.
\end{conjecture}

\begin{remark}
The conjecture is deliberately stated here without a claim about its complete
current status.  The archive contains several proved subfamilies, finite
checks, and stronger variants with additional parameters.  Those variants
should not be conflated: a proof for one dissection or one residue class does
not automatically settle the full three-polynomial assertion above.
\end{remark}

\section{How the modulus-specific material fits the general machinery}
\label{sec:extensions}

\begin{longtable}{>{\raggedright\arraybackslash}p{0.16\textwidth}
                  >{\raggedright\arraybackslash}p{0.34\textwidth}
                  >{\raggedright\arraybackslash}p{0.42\textwidth}}
\toprule
Archive theme & Structural source & Treatment in this article\\
\midrule
\endfirsthead
\toprule
Archive theme & Structural source & Treatment in this article\\
\midrule
\endhead
Finite $q$-products and Gaussian coefficients & Finite $q$-binomial theorem,
Cauchy augmentation, polynomial reciprocity & Fully proved in
\cref{thm:finiteqbt,lem:cauchyaug,prop:rectangle}.\\
\addlinespace
Basic hypergeometric summations & Heine transformation, $q$-Gauss,
$q$-Vandermonde, $q$-Pfaff--Saalsch\"utz & Fully proved in
\cref{thm:heine,cor:qgauss,cor:qvand1,prop:qvand2,thm:qps}.\\
\addlinespace
Theta and eta files & Jacobi products plus Poisson summation & Fully proved in
\cref{thm:jtp,thm:qpp,thm:thetaS,thm:eta-transform}.\\
\addlinespace
Rogers--Ramanujan and continued fractions & Unit Bailey pair, one chain step,
weak limit, continuants & Fully proved in
\cref{thm:RR,thm:RRCF}.\\
\addlinespace
Moduli $7,9,27$ & Odd-modulus Andrews--Gordon theorem & All members follow
from the proof of \cref{thm:AG}.\\
\addlinespace
Modulus $14$ and parity restrictions & Even-modulus Bressoud theory and
parity-sensitive Bailey transformations & Located as a distinct extension;
not silently identified with the odd-modulus theorem.\\
\addlinespace
Bilateral sums & Ramanujan's ${}_1\psi_1$ and theta specializations & Fully
proved in \cref{thm:1psi1}.\\
\addlinespace
Borwein-type signs & Cyclotomic residue dissections and reciprocal
polynomials & Structural constraints are proved in
\cref{prop:borweinreciprocity}; unproved positivity is labelled as
\cref{conj:borwein}.\\
\bottomrule
\end{longtable}

\subsection{Why modulus $14$ is genuinely different}
Odd-modulus Andrews--Gordon products arise from the bilateral exponent
\((2k+1)m^2+\lambda m\)/2 and an ordinary triple product.  Even-modulus
identities introduce a parity condition or a factor such as
$(-q;q)_N$ at the finite level.  Algebraically, this changes the Bailey
kernel; combinatorially, it changes the admissible frequencies of repeated
parts.  A correct modulus-$14$ treatment therefore needs a new lemma, not a
formal substitution $2k+1\mapsto14$ in \cref{thm:AG}.

\subsection{A reusable proof protocol}
For a new identity suggested by the archive, a reliable protocol is:
\begin{enumerate}[label=\arabic*.]
\item normalize every shifted factorial and every ${}_r\phi_s$;
\item search for a terminating polynomial form before taking an infinite
limit;
\item identify the exact Bailey pair and relative parameter;
\item verify the pair by inversion or a finite summation;
\item only then take weak limits and apply a product theorem;
\item separate coefficientwise convergence from numerical pointwise
convergence;
\item label sign claims as conjectural unless a coefficient-preserving
argument is supplied.
\end{enumerate}
This protocol is more important than any single formula: it turns a collection
of experimental identities into a dependency-controlled theory.

\section{Frontier problems suggested by the unified viewpoint}

The following problems are phrased so that proved statements and speculative
claims remain sharply separated.

\begin{problem}[Finite Andrews--Gordon certificates]\label{prob:finiteAG}
Construct, for every $(k,i)$, a finite polynomial identity whose
coefficientwise limit is \eqref{eq:AG}, whose summands are manifestly
nonnegative, and whose proof uses only Gaussian recurrences and a bounded
number of telescoping certificates.  Determine a canonical choice of
truncation for which the product side has a direct cylindric-partition or
lattice-path interpretation.
\end{problem}

\begin{problem}[An elementary even-modulus lowering step]\label{prob:evenlower}
Find a parity-sensitive analogue of \cref{lem:lowering} that turns the
odd-modulus Bailey-chain proof into a proof of the even-modulus Bressoud
family, with modulus $14$ as the first substantial test case.  The desired
operator should make the parity constraint visible at the level of the
Bailey kernel rather than introduce it only after summation.
\end{problem}

\begin{conjecture}[Sign-stable Borwein tails]\label{conj:signstable}
For each fixed offset $d\ge0$, the coefficients at distance $d$ from either
end of the reciprocal polynomials $A_n,B_n,C_n$ eventually stabilize as
$n\to\infty$, and the stabilized coefficients are nonnegative.
\end{conjecture}

\begin{remark}
The reciprocity proved in \cref{prop:borweinreciprocity} makes the two ends
equivalent.  The conjecture is weaker than full coefficientwise
nonnegativity, but it is suited to finite-product stabilization and can be
tested without extrapolating the central coefficients.
\end{remark}

\begin{problem}[Total positivity of Bailey kernels]\label{prob:TP}
For $0<q,a<1$, study minors of the lower-triangular matrix
\[
M_{n,r}(a,q)=\frac1{(q;q)_{n-r}(aq;q)_{n+r}}\qquad(0\le r\le n).
\]
Determine whether a natural diagonal normalization makes $M$ totally
positive.  A positive answer would turn many coefficient-sign statements
into consequences of variation-diminishing linear algebra; a negative answer
should identify the first obstructing minor and the correct weaker notion.
\end{problem}

\begin{problem}[Root-of-unity asymptotics]\label{prob:rootsunity}
Extend the exact modular acceleration \eqref{eq:qpoch-exact} to products and
Bailey sums as $q$ approaches a root of unity radially or tangentially.
Separate the dominant dilogarithmic exponential, the algebraic prefactor,
and the exponentially small resurgence sectors.  The finite forms in
\cref{prop:coefflimit} should serve as uniform error controls.
\end{problem}

\begin{problem}[Machine-checkable identity certificates]\label{prob:formal}
Formalize the dependency chain
\[
q\text{-binomial}\Rightarrow q\text{-Pfaff--Saalsch\"utz}
\Rightarrow\text{Bailey lemma}\Rightarrow\text{Andrews--Gordon}
\]
in a proof assistant.  The finite telescoping certificate in
\cref{lem:rogersfinite} is especially suitable: its verification reduces to
normalization of a rational function.  A useful library should keep formal
power series, analytic convergence, and meromorphic parameter continuation as
three distinct layers.
\end{problem}

\begin{problem}[Bilateral Bailey lattices]\label{prob:bilateral}
Combine the parameter-lowering idea with Ramanujan's ${}_1\psi_1$ summation to
construct bilateral Bailey chains whose limiting products have prescribed
residue classes and parity.  Identify when a bilateral chain collapses to a
unilateral Andrews--Gordon sum and when genuinely new negative-index
contributions survive.
\end{problem}

\section{Proof audit and dependency table}

\begin{longtable}{>{\raggedright\arraybackslash}p{0.28\textwidth}
                  >{\raggedright\arraybackslash}p{0.62\textwidth}}
\toprule
Result & Dependencies used in its proof\\
\midrule
\endfirsthead
\toprule
Result & Dependencies used in its proof\\
\midrule
\endhead
Finite $q$-binomial theorem & Pascal recurrence only.\\
Infinite $q$-binomial theorem & Coefficient recurrence and product
functional equation.\\
Heine transformation & Infinite $q$-binomial theorem plus an absolutely
convergent double-sum interchange.\\
$q$-Gauss and $q$-Vandermonde & Heine transformation, product splitting, and
finite reversal.\\
Rogers finite product lemma & Pascal recurrence and an explicit rational
summand certificate.\\
$q$-Pfaff--Saalsch\"utz & Rogers finite product lemma after denominator
clearing.\\
Jacobi triple product & Euler expansions, a Laurent functional equation, and
$q$-Gauss for the constant term.\\
Quintuple product & Two triple products, a three-step Laurent recurrence, and
Euler's pentagonal theorem.\\
Bailey inversion & A finite $q$-difference annihilation lemma.\\
Bailey lemma & $q$-Pfaff--Saalsch\"utz.\\
Rogers--Ramanujan & Unit Bailey pair, one chain step, weak Bailey lemma, and
Jacobi triple product.\\
Andrews--Gordon & Unit pair relative to $q$, chain steps, parameter lowering,
weak Bailey lemma, and triple product.\\
Theta inversion & Poisson summation and the Gaussian Fourier transform.\\
Eta inversion and $(q;q)_\infty$ asymptotics & Theta inversion and Jacobi
product formulas.\\
Ramanujan ${}_1\psi_1$ & $q$-binomial theorem at a sequence of parameter
values and analytic continuation.\\
\bottomrule
\end{longtable}

\section{Conclusion}

The central lesson of the archive is not merely that many striking products
exist.  It is that a small number of finite mechanisms repeatedly generate
them.  The finite $q$-binomial theorem supplies the elementary algebra;
Heine's transformation and the finite Rogers lemma supply the decisive
summations; Bailey inversion turns a kernel relation into a reusable
transformation; and Jacobi's product converts bilateral quadratic exponents
into residue-class products.  Poisson summation then connects the same
products to modular transformations and sharp asymptotics.

The parameter-lowering lemma makes the full Andrews--Gordon family especially
transparent: the position at which one changes from a Bailey pair relative to
$q$ to one relative to $1$ is exactly the position at which the linear terms
begin in the multisum.  This unifies the modulus $5,7,9,$ and $27$ material.
At the same time, it explains why even modulus and parity-restricted files
need genuinely different machinery.  Finally, the separation of theorems
from conjectural sign patterns provides a safe foundation for further
symbolic experimentation and formal verification.

\appendix
\section{Analytic details used repeatedly}

\begin{lemma}[Normal convergence of shifted-factorial series]\label{lem:normal}
Let $|q|<1$.  On compact subsets avoiding denominator zeros, every terminating
basic hypergeometric series is rational in its parameters, and every
nonterminating series used in this article converges normally in the stated
open domain.
\end{lemma}
\begin{proof}
Termination is immediate.  For ${}_2\phi_1$, the ratio of successive terms
tends to $z$, so compact subsets of $|z|<1$ admit a uniform ratio bound below
$1$.  For bilateral ${}_1\psi_1$, the positive and negative tails have limiting
ratios $z$ and $b/(az)$ after reversal; compact subannuli of
$|b/a|<|z|<1$ again admit uniform bounds.  Products converge normally by
\cref{lem:prodconv}.  These estimates also justify every sum interchange and
parameter limit made above.
\end{proof}

\begin{lemma}[Removable singularity principle]\label{lem:removable}
Suppose a finite rational identity holds for parameters in a nonempty open
set and both sides become polynomials after multiplication by the same
explicit denominator.  Then the denominator-cleared identity holds for all
parameter values, and any value at a zero of the denominator is determined by
cancellation or a limit.
\end{lemma}
\begin{proof}
The difference of the denominator-cleared sides is a polynomial vanishing on
an open set, hence is identically zero.  Dividing back gives equality wherever
the denominator is nonzero; cancelled values are removable by continuity.
\end{proof}

\section{Notation index}

\begin{center}
\begin{tabularx}{0.92\textwidth}{@{}lX@{}}
\toprule
Symbol & Meaning\\
\midrule
$(a;q)_n$ & finite $q$-shifted factorial
$\prod_{j=0}^{n-1}(1-aq^j)$\\
$(a;q)_\infty$ & convergent infinite shifted factorial for $|q|<1$\\
$(a_1,\ldots,a_r;q)_n$ & product of shifted factorials with a common base\\
$\GaussianBinomial{n}{k}{q}$ & Gaussian or $q$-binomial coefficient\\
${}_r\phi_s$ & unilateral basic hypergeometric series in the normalization of
\cref{thm:heine}\\
${}_1\psi_1$ & bilateral basic hypergeometric series\\
$j(z;q)$ & Jacobi product $(q,z,q/z;q)_\infty$\\
$\vartheta_2,\vartheta_3,\vartheta_4$ & Jacobi theta constants with nome
$\EulerE^{\pi i\tau}$\\
$\eta(\tau)$ & Dedekind eta function\\
$(\alpha,\beta)$ relative to $a$ & Bailey pair satisfying
\eqref{eq:baileypair}\\
$N_j$ & tail sum $n_j+\cdots+n_{k-1}$ in Andrews--Gordon\\
\bottomrule
\end{tabularx}
\end{center}

\begin{thebibliography}{99}

\bibitem{Andrews1974}
G.~E. Andrews,
\newblock An analytic generalization of the Rogers--Ramanujan identities for
odd moduli,
\newblock \emph{Proceedings of the National Academy of Sciences USA}
\textbf{71} (1974), 4082--4085.

\bibitem{AndrewsPartitions}
G.~E. Andrews,
\newblock \emph{The Theory of Partitions},
\newblock Encyclopedia of Mathematics and its Applications, Vol.~2,
Addison--Wesley, 1976; Cambridge University Press reprint, 1998.

\bibitem{Bailey1949}
W.~N. Bailey,
\newblock Identities of the Rogers--Ramanujan type,
\newblock \emph{Proceedings of the London Mathematical Society} (2)
\textbf{50} (1949), 1--10.

\bibitem{BorweinRestricted}
P.~B. Borwein,
\newblock Some restricted partition functions,
\newblock \emph{Journal of Number Theory} \textbf{35} (1990), 228--240.

\bibitem{Bressoud1979}
D.~M. Bressoud,
\newblock A generalization of the Rogers--Ramanujan identities for all
moduli,
\newblock \emph{Journal of Combinatorial Theory, Series A} \textbf{27}
(1979), 64--68.

\bibitem{BressoudMemoir}
D.~M. Bressoud,
\newblock Analytic and combinatorial generalizations of the
Rogers--Ramanujan identities,
\newblock \emph{Memoirs of the American Mathematical Society} \textbf{24},
no.~227 (1980).

\bibitem{GasperRahman}
G.~Gasper and M.~Rahman,
\newblock \emph{Basic Hypergeometric Series}, second edition,
\newblock Encyclopedia of Mathematics and its Applications, Vol.~96,
Cambridge University Press, 2004.

\bibitem{Gordon1961}
B.~Gordon,
\newblock A combinatorial generalization of the Rogers--Ramanujan identities,
\newblock \emph{American Journal of Mathematics} \textbf{83} (1961),
393--399.

\bibitem{Heine1847}
E.~Heine,
\newblock Untersuchungen \"uber die Reihe
$1+\frac{(1-q^\alpha)(1-q^\beta)}{(1-q)(1-q^\gamma)}x+\cdots$,
\newblock \emph{Journal f\"ur die reine und angewandte Mathematik}
\textbf{34} (1847), 285--328.

\bibitem{Jacobi1829}
C.~G.~J. Jacobi,
\newblock \emph{Fundamenta Nova Theoriae Functionum Ellipticarum},
\newblock K\"onigsberg, 1829.

\bibitem{RamanujanCollected}
S.~Ramanujan,
\newblock \emph{Collected Papers},
\newblock edited by G.~H. Hardy, P.~V. Seshu Aiyar, and B.~M. Wilson,
Cambridge University Press, 1927.

\bibitem{Rogers1894}
L.~J. Rogers,
\newblock Second memoir on the expansion of certain infinite products,
\newblock \emph{Proceedings of the London Mathematical Society} (1)
\textbf{25} (1894), 318--343.

\bibitem{Slater1952}
L.~J. Slater,
\newblock Further identities of the Rogers--Ramanujan type,
\newblock \emph{Proceedings of the London Mathematical Society} (2)
\textbf{54} (1952), 147--167.

\end{thebibliography}

\end{document}
